\chapter{Factoring Polynomials}


\section{Factoring Polynomials}
\textbf{Factoring} is when you rewrite a polynomial as a product of polynomials.\\

\subsection{Greatest Common Factor}
The first method is called \textbf{greatest common factor} as you factor out the largest number in the polynomial.\\
\\
Example: $5x + 5y = 5(x + y)$\\

\subsection{Difference of Squares}
The product of conjugates produces a pattern called a \textbf{difference of squares}.\\
\\
Example: $(a + b)(a - b) = a^2 - b^2$\\
\\

\subsection{Sum or Difference of Cubes}
\textbf{Sum of Cubes} is the form $(a)^3 + (b)^3$.\\
\textbf{Difference of Cubes} is the form of $(a)^3 - (b)^3$\\
\\
They have similar factored patterns where they follow the same rules.\\
The first binomial has the same sign, but the trinomial has the \textit{opposite} sign.\\
$(a)^3 + (b)^3 = [(a) + (b)][(a)^2 - (a)(b) + (b)^2]$\\
\\
$(a)^3 - (b)^3 = [(a) + (b)][(a)^2 + (a)(b) + (b)^2]$\\


\textbf{Trinomials of the Form $x^2 + bx + c$}
Example: $2x^2 + 16x + 24$ \\
\\
Step 1: See if you can take out the Greatest Common Factor to make the problem easier. In this case you can take out a \textit{2} from all three numbers.

To factor these you first distribute the first $x^2$ into two parenthesis to start:\\ Don't forget to put the 2 you factored out in the beginning. So far we have $2(x^2 + 8x + 12)$.
	$2(x )(x )$\\
\\
Now you think of two numbers which can add up the middle number but also multiply to the last number. Keep in mind you can switch up the symbols to make either one positive or negative. Some factors to try are:
$(1)(12)$, $(-1)(-12)$, $(2)(6)$, $(3)(4)$\\
\\
Plugging in 2 and 6 we can see that they add up to 8 and multiply to 12 giving us our factored trinomial, which keep in mind is just two binomial.\\
\\
$2(x + 2)(x + 6)$\\
\\
\subsection{Sqaure Trinomials}
Recall $(x + y)^2 = (x)^2 + 2(x)(y) + (y)^2$\\
\\
Therefore, it is easy to see if an equation is a square trinomial or binomial.\\
\\
Example: $4x^2 - 20x + 25$ can be deduced down to squares: $4x^2 = (2x)^2$ \& $25 = (-5)^2$\\
Leaving us with the final step of putting them in the form $2xy$ to see if it fits: 2(2x)(-5) which does equal -20x like the original equation. This means $(2x - 5)^2$ is the square binomial to our original trinomial.\\
\\

\subsection{Factoring by Regrouping}
There are two main things to look for in attempting to factor a polynomial of four or more terms with no GCF. It can be regrouped into two groups of two terms or two groups where one has three terms aka a trinomial. \\
\\
Example: (az + ay) + (bz + by)\\
\\
We can factor out both an 'a' and a 'b'. Leaving us with $a(y + z) + b(y + z)$. Now we see we can see that we can factor out (y + z) leaving us now with $(y + z)(a + b)$. \\

